% Part: first-order-logic
% Chapter: axiomatic-deduction
% Section: proving-things

\documentclass[../../../include/open-logic-section]{subfiles}

\begin{document}

\iftag{FOL}
      {\olfileid{fol}{axd}{pro}}
      {\olfileid{pl}{axd}{pro}}

\olsection{\usetoken{P}{derivation} دیاں مثالاں}


\begin{ex}
فرض کرو اسیں $(\lnot !D \lor !E) \lif (!D \lif
!E)$ ثابت کرنا چاہندے آں۔ ظاہر اے ایہ ساڈے کسے وی مسلّمے دی مثال نہیں؛
ایس لئی ایہنوں !!{derive} کرن لئی \MP{} دا قاعدہ ورتنا پئے گا۔ ساڈا
اکّو قاعدہ~MP اے، جیہڑا $!A$ تے $!A \lif !B$ دتے ہون اُتے~$!B$
نوں درست قدم ٹھہران دی اجازت دیندا اے۔ اک طریقہ ایہ ہو سکدا اے کہ
\olref[prp]{ax:lor3} وچ $!A$ دی تھاں $\lnot
!D$، $!B$ دی تھاں $!E$ تے $!C$ دی تھاں $!D \lif !E$ رکھیے،
یعنی ایہ مثال:
\[
(\lnot !D \lif (!D \lif !E)) \lif ((!E \lif (!D \lif !E)) \lif ((\lnot
!D \lor !E) \lif (!D \lif !E))).
\]
کیوں؟ MP نوں دو واری ورتن نال آخری حصہ ملدا اے، تے اوہی اسیں
چاہندے آں۔ اسیں سوکھ نال ویکھ سکدے آں کہ $\lnot !D \lif (!D \lif !E)$،
\olref[prp]{ax:lnot2} دی مثال اے، تے $!E \lif (!D \lif !E)$،
\olref[prp]{ax:lif1} دی مثال اے۔ ایس لئی ساڈا !!{derivation} ایہ اے:
\begin{derivation}
  1. &  $\lnot !D \lif (!D \lif !E)$ & \olref[prp]{ax:lnot2} \\
  2. & $(\lnot !D \lif (!D \lif !E)) \lif {}$\\
  &\qquad $((!E \lif (!D \lif !E)) \lif ((\lnot !D \lor !E) \lif (!D \lif !E)))$ & \olref[prp]{ax:lor3}\\
  3. & $(!E \lif (!D \lif !E)) \lif ((\lnot !D \lor !E) \lif (!D \lif !E))$ & 1, 2, \MP\\
  4. & $!E \lif (!D \lif !E)$ & \olref[prp]{ax:lif1}\\
  5. & $(\lnot !D \lor !E) \lif (!D \lif !E)$ & 3, 4, \MP
\end{derivation}
\end{ex}

\begin{ex}\ollabel{ex:identity}
آؤ $!D \lif !D$ دا !!a{derivation} لبھن دی کوشش کریے۔ ایہ کسے
مسلّمے دی مثال نہیں، ایس لئی ایہنوں !!{derive} کرن لئی \MP{} ورتنا
پئے گا۔ \olref[prp]{ax:lif1}،~$!A \lif !B$ دے روپ دا اک مسلّمہ
اے جیہدے اُتے اسیں~\MP لا سکدے آں۔ پر ایس حالت وچ \MP{} جیہڑے
$!B$ نوں درست قدم ٹھہرائے، اوہ لازماً~$!D \lif !D$ ہونا چاہیدا
اے، کیوں جے اسیں ایہو !!{derive} کرنا چاہندے آں۔ ایس دا مطلب اے
کہ $!A$ وی $!D$ ہونا چاہیدا اے؛ یعنی اسیں
\olref[prp]{ax:lif1} دی ایہ مثال ویکھ سکدے آں:
\[
!D \lif (!D \lif !D)
\]
\MP ورتن لئی ساڈے کول اوہدے مطابق دوجے مقدمے، یعنی~$!A$، دا
جواز وی ہونا چاہیدا اے۔ پر ایتھے اوہ~$!D$ ہووے گا، تے اسیں
اکیلا~$!D$ !!{derive} نہیں کر سکاں گے۔ ایس لئی سانوں ہور
طریقہ چاہیدا اے۔

صرف~$\lif$ نال متعلق دوجا مسلّمہ \olref[prp]{ax:lif2} اے، یعنی
\[
(!A \lif (!B \lif !C)) \lif ((!A \lif !B) \lif (!A \lif !C))
\]
\MP{} نوں دو واری ورت کے اسیں آخری اندرلا شرطیہ حاصل کر سکدے آں۔
فیر ایہدا مطلب اے کہ سانوں \olref[prp]{ax:lif2} دی ایسی مثال
چاہیدی اے جیہدے وچ $!A \lif !C$، اوہ !!{formula}~$!D \lif !D$
ہووے جیہنوں اسیں ثابت کرنا چاہندے آں۔ ایس لئی $!A$ تے $!C$
دوویں~$!D$ نیں۔ ہُن $!B$ کیہڑا چُنیے تاں جو $!A \lif (!B \lif !C)$
تے $!A \lif !B$، یعنی ایتھے $!D \lif (!B \lif !D)$ تے $!D \lif !B$،
دوویں !!{derivable} ہون؟ پہلا تاں $!B$ دی ہر چُون لئی
\olref[prp]{ax:lif1} دی مثال اے۔ تے $!D \lif !B$ وی
\olref[prp]{ax:lif1} دی مثال ہووے گا جے $!B$، $(!D \lif !D)$ ہووے۔
ایس لئی ساڈا !!{derivation} ایہ اے:
\begin{derivation}
  1. & $!D \lif ((!D \lif !D) \lif !D)$ & \olref[prp]{ax:lif1}\\
  2. & $(!D \lif ((!D \lif !D) \lif !D)) \lif {}$\\
  & \qquad $((!D \lif (!D \lif !D)) \lif (!D \lif !D))$ & \olref[prp]{ax:lif2}\\
  3. & $(!D \lif (!D \lif !D)) \lif (!D \lif !D)$ & 1, 2, \MP\\
  4. & $!D \lif (!D \lif !D)$ & \olref[prp]{ax:lif1}\\
  5. & $!D \lif !D$ & 3, 4, \MP
\end{derivation}
\end{ex}

\begin{ex}\ollabel{ex:chain}
کدے اسیں دکھانا چاہندے آں کہ کجھ ہور !!{formula}s~$\Gamma$ توں
کسے !!{formula} دا !!a{derivation} موجود اے۔ مثال لئی، آؤ
دکھائیے کہ اسیں $!A \lif !C$، $\Gamma = \{!A \lif !B,
!B \lif !C\}$ توں !!{derive} کر سکدے آں۔
\begin{derivation}
  1. & $!A \lif !B$ & \Hyp\\
  2. & $!B \lif !C$ & \Hyp\\
  3. & $(!B \lif !C) \lif (!A \lif (!B \lif !C))$ & \olref[prp]{ax:lif1} \\
  4. & $!A \lif (!B \lif !C)$ & 2, 3, \MP\\
  5. & $(!A \lif (!B \lif !C)) \lif {}$\\
  & \qquad $((!A \lif !B) \lif (!A \lif !C))$ & \olref[prp]{ax:lif2}\\
  6. & $((!A \lif !B) \lif (!A \lif !C))$ & 4, 5, \MP\\
  7. & $!A \lif !C$ & 1, 6, \MP
\end{derivation}
``\Hyp'' نال نشان لگیاں سطراں (``مفروضہ'' لئی) دسدیاں
نیں کہ اس سطر دا !!{formula}،~$\Gamma$ دا !!a{element} اے۔
\end{ex}

\begin{prop}
  \ollabel{prop:chain} جے $\Gamma \Proves !A \lif !B$ تے $\Gamma
  \Proves !B \lif !C$، تاں $\Gamma \Proves !A \lif !C$
\end{prop}

\begin{proof}
  فرض کرو $\Gamma \Proves !A \lif !B$ تے $\Gamma \Proves !B \lif
  !C$۔ تاں $!A \lif !B$ دا !!a{derivation}~$\Gamma$ توں موجود
  اے، تے $!B \lif !C$ دا !!a{derivation}~$\Gamma$ توں وی۔ اوہناں
  نوں اک دوجے دے پچھے جوڑ کے اکّو !!{derivation} بنا لو۔ ہُن پچھلی
  مثال دے !!{derivation} دیاں سطراں 3--7 شامل کرو۔ ایس نال
  $!A \lif !C$ دا !!a{derivation}، جیہدے وچ ایہ فارمولا نویں
  !!{derivation} دی آخری سطر اے،~$\Gamma$ توں ملدا اے۔ دھیان دیو
  کہ سطراں 4 تے~7 دے جواز درست رہندے نیں جے سطر نمبر~1 دے
  حوالے نوں $!A \lif !B$ دے !!{derivation} دی آخری سطر دے
  حوالے نال، تے سطر نمبر~2 دے حوالے نوں $!B \lif !C$ دے
  !!{derivation} دی آخری سطر دے حوالے نال بدل دتا جائے۔
\end{proof}

\begin{prob}
مسلّماں توں !!{derivation}s پیش کر کے دکھاؤ کہ ہِٹھ دتیاں گلاں
درست نیں:
\begin{enumerate}
\item $(!A \land !B) \lif (!B \land !A)$
\item $((!A \land !B) \lif !C) \lif (!A \lif (!B \lif !C))$
\item $\lnot(!A \lor !B) \lif \lnot !A$
\end{enumerate}
\end{prob}



\end{document}
